% ══════════════════════════════════════════════════════════════════════════════
%                         فصل هشتم
%                 چپ نو و جنبش‌های معاصر
% ══════════════════════════════════════════════════════════════════════════════

\chapter{چپ نو و جنبش‌های معاصر}
\label{ch:new-left}

\begin{flushright}
\textit{«زیر سنگفرش، ساحل است.»}

— شعار مه ۱۹۶۸ پاریس
\end{flushright}

% ══════════════════════════════════════════════════════════════════════════════
%                         مقدمه فصل
% ══════════════════════════════════════════════════════════════════════════════

\lettrine[lines=3, lraise=0.1, nindent=0.5em]{\textcolor{PurpleAccent}{د}}{ر} 
دهه ۱۹۶۰، چپ جدیدی متولد شد که نه با چپ سنتی کمونیستی راحت بود، نه با سوسیال‌دموکراسی اداری و خسته. این «چپ نو» از دانشگاه‌ها و خیابان‌ها برآمد، زبان جدیدی داشت، و پرسش‌های جدیدی می‌پرسید: آزادی جنسی، حقوق زنان، محیط زیست، نژادپرستی، امپریالیسم. این فصل داستان تولد، اوج و تحول چپ نو را روایت می‌کند—از مه ۱۹۶۸ پاریس تا جنبش اشغال وال‌استریت.

% ══════════════════════════════════════════════════════════════════════════════
\section{شکست چپ سنتی و تولد چپ نو}
% ══════════════════════════════════════════════════════════════════════════════

\subsection{چرا چپ نو؟}

در دهه ۱۹۶۰، نسل جدیدی از فعالان با چپ موجود مشکل داشتند:

\begin{table}[htbp]
\centering
\caption{چپ سنتی در برابر چپ نو}
\label{tab:old-vs-new-left}
\begin{tabular}{r|c|c}
\toprule
\textbf{موضوع} & \textbf{چپ سنتی} & \textbf{چپ نو} \\
\midrule
محور اصلی & طبقه کارگر & دانشجویان، جوانان، حاشیه‌ها \\
هدف & تصرف قدرت دولتی & تغییر فرهنگ و زندگی روزمره \\
سازمان‌دهی & حزب متمرکز & شبکه‌ای، افقی، خودجوش \\
الگو & شوروی یا سوسیال‌دموکراسی & کوبا، چین، جهان سوم \\
موضوعات & اقتصاد، کار، دستمزد & جنسیت، نژاد، محیط زیست، جنگ \\
سبک & رسمی، جدی، بوروکراتیک & شاد، خلاق، ضدفرهنگ \\
\bottomrule
\end{tabular}
\end{table}

\begin{defbox}[تعریف: چپ نو]
جنبشی که از دهه ۱۹۵۰ در غرب شکل گرفت و ویژگی‌هایش عبارت بود از:
\begin{itemize}[nosep]
    \item نقد هم سرمایه‌داری و هم «سوسیالیسم واقعاً موجود» (شوروی)
    \item تمرکز بر \textbf{آزادی فردی} و \textbf{خودمختاری}
    \item توجه به مسائل فرهنگی، نه صرفاً اقتصادی
    \item مبارزه علیه نژادپرستی، جنسیت‌گرایی، امپریالیسم
    \item روش‌های نوین: اقدام مستقیم، نافرمانی مدنی، ضدفرهنگ
\end{itemize}
\end{defbox}

% ══════════════════════════════════════════════════════════════════════════════
\section{ریشه‌های فکری چپ نو}
% ══════════════════════════════════════════════════════════════════════════════

\subsection{مکتب فرانکفورت}

\begin{historybox}[مکتب فرانکفورت]
گروهی از متفکران مارکسیست آلمانی که «نظریه انتقادی» را بنیان نهادند:

\textbf{مؤسسه تحقیقات اجتماعی (۱۹۲۳):} فرانکفورت → تبعید به آمریکا → بازگشت پس از جنگ

\textbf{چهره‌های کلیدی:}
\begin{itemize}[nosep]
    \item ماکس هورکهایمر
    \item تئودور آدورنو
    \item هربرت مارکوزه
    \item اریش فروم
    \item والتر بنیامین
    \item (نسل دوم:) یورگن هابرماس
\end{itemize}

\textbf{ایده محوری:} مارکسیسم باید با روان‌کاوی، جامعه‌شناسی و زیبایی‌شناسی ترکیب شود تا بتواند سرمایه‌داری متأخر را نقد کند.
\end{historybox}

\begin{personbox}[هربرت مارکوزه (۱۸۹۸-۱۹۷۹)]

\textbf{ملیت:} آلمانی-آمریکایی

\textbf{آثار کلیدی:} \book{اروس و تمدن} (۱۹۵۵)، \book{انسان تک‌ساحتی} (۱۹۶۴)

\textbf{لقب:} «پدر چپ نو»

\mediumspace

\person{مارکوزه} تأثیرگذارترین فیلسوف بر جنبش‌های ۱۹۶۸ بود:

\begin{itemize}[nosep]
    \item \textbf{انسان تک‌ساحتی:} سرمایه‌داری متأخر انسان‌ها را به مصرف‌کننده تبدیل کرده
    \item \textbf{سرکوب مازاد:} جامعه صنعتی بیش از حد لازم سرکوب می‌کند
    \item \textbf{تحمل سرکوبگرانه:} تحمل ظاهری که نقد رادیکال را بی‌اثر می‌کند
    \item \textbf{امتناع بزرگ:} نه گفتن به کل نظام، نه فقط اصلاح آن
    \item \textbf{نیروهای انقلابی جدید:} دانشجویان، سیاهان، زنان، جهان سوم
\end{itemize}

\mediumspace

\textbf{نقل‌قول:} \textit{«آزاد بودن در واقع یعنی آزاد بودن از فقر، آزاد بودن از نیاز به تلاش برای بقا.»}
\end{personbox}

\begin{quotebox}[هربرت مارکوزه، انسان تک‌ساحتی]
\textit{«تحت حکومت یک کل سرکوبگر، آزادی می‌تواند به ابزار قدرتمند سلطه تبدیل شود. دامنه انتخاب‌های باز به فرد، معیار تعیین‌کننده آزادی نیست، بلکه آنچه می‌تواند انتخاب شود و آنچه واقعاً توسط فرد انتخاب می‌شود، معیار است.»}
\end{quotebox}

\subsection{نقد تمدن صنعتی}

\begin{defbox}[تعریف: نقد تمدن صنعتی]
مکتب فرانکفورت و چپ نو استدلال کردند:
\begin{itemize}[nosep]
    \item مشکل فقط سرمایه‌داری نیست—\textbf{تمدن صنعتی} خود سرکوبگر است
    \item شوروی هم صنعتی بود، پس سرکوبگر
    \item عقلانیت ابزاری (تسلط بر طبیعت و انسان) مشکل اصلی است
    \item فرهنگ توده‌ای (صنعت فرهنگ) آگاهی را کنترل می‌کند
    \item رهایی، نه فقط اقتصادی، بلکه فرهنگی و روان‌شناختی است
\end{itemize}
\end{defbox}

% ══════════════════════════════════════════════════════════════════════════════
\section{مه ۱۹۶۸: لحظه انقلابی}
% ══════════════════════════════════════════════════════════════════════════════

\begin{figure}[htbp]
\centering
\begin{tikzpicture}[scale=0.85]

% عنوان
\node[font=\large\bfseries, text=DarkGray] at (0, 8) {\fa{۱۹۶۸ در جهان: سال انقلاب جهانی}};

% نقشه نمادین
% پاریس
\fill[LeftRed] (-5, 5) circle (12pt);
\node[text=white, font=\scriptsize\bfseries] at (-5, 5) {\fa{پاریس}};
\node[below, text=DarkGray, font=\tiny, text width=2cm, align=center] at (-5, 4.3) {
    \fa{مه ۶۸}\\
    \fa{دانشجویان + کارگران}
};

% برلین
\fill[LeftRed] (-2.5, 6) circle (10pt);
\node[text=white, font=\scriptsize\bfseries] at (-2.5, 6) {\fa{برلین}};
\node[below, text=DarkGray, font=\tiny] at (-2.5, 5.4) {\fa{جنبش دانشجویی}};

% پراگ
\fill[CenterGreen] (0, 5.5) circle (10pt);
\node[text=white, font=\scriptsize\bfseries] at (0, 5.5) {\fa{پراگ}};
\node[below, text=DarkGray, font=\tiny, text width=2cm, align=center] at (0, 4.8) {
    \fa{بهار پراگ}\\
    \fa{سرکوب شوروی}
};

% مکزیکو
\fill[LeftRedDark] (-4, 2.5) circle (10pt);
\node[text=white, font=\scriptsize\bfseries] at (-4, 2.5) {\fa{مکزیکو}};
\node[below, text=DarkGray, font=\tiny, text width=2cm, align=center] at (-4, 1.8) {
    \fa{تلاتلولکو}\\
    \fa{کشتار دانشجویان}
};

% آمریکا
\fill[LeftRed] (3, 4) circle (12pt);
\node[text=white, font=\scriptsize\bfseries] at (3, 4) {\fa{آمریکا}};
\node[below, text=DarkGray, font=\tiny, text width=2.5cm, align=center] at (3, 3.2) {
    \fa{ضد جنگ ویتنام}\\
    \fa{حقوق مدنی}\\
    \fa{ترور کینگ و کندی}
};

% ژاپن
\fill[LeftRedLight] (5.5, 5.5) circle (10pt);
\node[text=white, font=\scriptsize\bfseries] at (5.5, 5.5) {\fa{توکیو}};
\node[below, text=DarkGray, font=\tiny] at (5.5, 4.9) {\fa{زنگاکورن}};

% شعارها
\node[text=PurpleAccent, font=\small\bfseries, text width=10cm, align=center] at (0, 0.8) {
    \fa{«واقع‌بین باش، ناممکن را بخواه!»}\\
    \fa{«ممنوع کردن ممنوع است!»}\\
    \fa{«زیر سنگفرش، ساحل است.»}
};

\end{tikzpicture}
\caption{جغرافیای جنبش‌های ۱۹۶۸}
\label{fig:1968-map}
\end{figure}

\begin{historybox}[مه ۱۹۶۸ پاریس]
\textbf{آغاز:} اعتراض دانشجویان دانشگاه نانتر به قوانین خوابگاه

\textbf{گسترش:}
\begin{itemize}[nosep]
    \item ۳ مه: پلیس دانشگاه سوربن را می‌بندد
    \item ۱۰ مه: شب سنگرها—دانشجویان در کوارتیه لاتن
    \item ۱۳ مه: اعتصاب عمومی—۱۰ میلیون کارگر
    \item ۲۷ مه: توافق گرونل—افزایش دستمزد ۳۵٪
    \item ۳۰ مه: تظاهرات حمایت از دوگل—پایان بحران
\end{itemize}

\textbf{نتیجه:} انقلاب «شکست» خورد، اما فرهنگ را تغییر داد.

\textbf{چهره‌ها:} دانیل کوهن-بندیت («دنی سرخ»)، جی‌پی سارتر (حمایت)
\end{historybox}

\begin{debatebox}[مه ۶۸: انقلاب یا سایکودرام؟]
\textbf{ارزیابی مثبت:}
\begin{itemize}[nosep]
    \item نشان داد انقلاب در غرب ممکن است
    \item فرهنگ را دگرگون کرد: آزادی جنسی، فمینیسم، اکولوژی
    \item نسلی از فعالان را تربیت کرد
    \item پایان «سال‌های شکوهمند» سرمایه‌داری بی‌چون‌وچرا
\end{itemize}

\textbf{ارزیابی منفی:}
\begin{itemize}[nosep]
    \item «انقلابی» که قدرت نگرفت، انقلاب نیست
    \item شورش فرزندان بورژوازی بود، نه انقلاب طبقاتی
    \item نتیجه‌اش نئولیبرالیسم و فردگرایی شد (نقد از چپ)
    \item «سوسیالیسم شامپاین» (نقد از راست)
\end{itemize}
\end{debatebox}

\subsection{جنبش ضد جنگ ویتنام}

\begin{historybox}[جنبش ضد جنگ ویتنام در آمریکا]
\begin{itemize}[nosep]
    \item \textbf{SDS:} دانشجویان برای جامعه دموکراتیک
    \item \textbf{بیانیه پورت هورون (۱۹۶۲):} مانیفست چپ نو آمریکایی
    \item \textbf{تظاهرات:} صدها هزار نفر در واشنگتن
    \item \textbf{مقاومت:} فرار از سربازی، سوزاندن کارت نظام‌وظیفه
    \item \textbf{کنت استیت (۱۹۷۰):} گارد ملی ۴ دانشجو را کشت
\end{itemize}

\textbf{نتیجه:} کمک به پایان جنگ، اما جنبش پراکنده شد.
\end{historybox}

% ══════════════════════════════════════════════════════════════════════════════
\section{جنبش‌های هویتی}
% ══════════════════════════════════════════════════════════════════════════════

چپ نو توجه را از طبقه به \keyword{هویت} چرخاند: نژاد، جنسیت، گرایش جنسی.

\subsection{جنبش حقوق مدنی سیاه‌پوستان}

\begin{figure}[htbp]
\centering
\begin{tikzpicture}[scale=0.85]

% عنوان
\node[font=\large\bfseries, text=DarkGray] at (0, 7) {\fa{تایم‌لاین: جنبش حقوق مدنی آمریکا}};

% خط زمان
\draw[line width=4pt, MediumGray] (-7, 5) -- (7, 5);

% رویدادها
% براون
\fill[CenterGreen] (-6, 5) circle (9pt);
\node[above, text=CenterGreen, font=\scriptsize\bfseries] at (-6, 5.4) {\fa{۱۹۵۴}};
\node[below, text width=2cm, align=center, text=DarkGray, font=\tiny] at (-6, 4.4) {
    \fa{حکم براون}\\
    \fa{پایان تفکیک مدارس}
};

% مونتگمری
\fill[CenterGreenLight] (-4, 5) circle (9pt);
\node[above, text=CenterGreenLight, font=\scriptsize\bfseries] at (-4, 5.4) {\fa{۱۹۵۵}};
\node[below, text width=2cm, align=center, text=DarkGray, font=\tiny] at (-4, 4.4) {
    \fa{بایکوت اتوبوس}\\
    \fa{رزا پارکس}
};

% مارش واشنگتن
\fill[PurpleAccent] (-1.5, 5) circle (10pt);
\node[above, text=PurpleAccent, font=\scriptsize\bfseries] at (-1.5, 5.5) {\fa{۱۹۶۳}};
\node[below, text width=2.2cm, align=center, text=DarkGray, font=\tiny] at (-1.5, 4.3) {
    \fa{مارش واشنگتن}\\
    \fa{«رویایی دارم»}
};

% قانون حقوق مدنی
\fill[CenterGreen] (1, 5) circle (9pt);
\node[above, text=CenterGreen, font=\scriptsize\bfseries] at (1, 5.4) {\fa{۱۹۶۴}};
\node[below, text width=2cm, align=center, text=DarkGray, font=\tiny] at (1, 4.4) {
    \fa{قانون}\\
    \fa{حقوق مدنی}
};

% مالکوم ایکس
\fill[LeftRedDark] (3, 5) circle (9pt);
\node[above, text=LeftRedDark, font=\scriptsize\bfseries] at (3, 5.4) {\fa{۱۹۶۵}};
\node[below, text width=2cm, align=center, text=DarkGray, font=\tiny] at (3, 4.4) {
    \fa{ترور}\\
    \fa{مالکوم ایکس}
};

% پلنگ سیاه
\fill[LeftRedDark] (5, 5) circle (10pt);
\node[above, text=LeftRedDark, font=\scriptsize\bfseries] at (5, 5.5) {\fa{۱۹۶۶}};
\node[below, text width=2cm, align=center, text=DarkGray, font=\tiny] at (5, 4.3) {
    \fa{پلنگان سیاه}\\
    \fa{قدرت سیاه}
};

% کینگ
\fill[MediumGray] (6.5, 5) circle (9pt);
\node[above, text=MediumGray, font=\scriptsize\bfseries] at (6.5, 5.4) {\fa{۱۹۶۸}};
\node[below, text width=2cm, align=center, text=DarkGray, font=\tiny] at (6.5, 4.4) {
    \fa{ترور}\\
    \fa{مارتین لوتر کینگ}
};

\end{tikzpicture}
\caption{سیر جنبش حقوق مدنی سیاهان آمریکا}
\label{fig:civil-rights-timeline}
\end{figure}

\begin{personbox}[مارتین لوتر کینگ جونیور (۱۹۲۹-۱۹۶۸)]

\textbf{ملیت:} آمریکایی

\textbf{روش:} نافرمانی مدنی مسالمت‌آمیز (الهام از گاندی)

\textbf{جایزه صلح نوبل:} ۱۹۶۴

\mediumspace

\textbf{دستاوردها:}
\begin{itemize}[nosep]
    \item رهبری بایکوت مونتگمری
    \item مارش واشنگتن و سخنرانی «رویایی دارم»
    \item فشار برای قانون حقوق مدنی ۱۹۶۴
    \item در سال‌های آخر: نقد جنگ ویتنام و فقر (رادیکال‌تر شد)
\end{itemize}

\mediumspace

\textbf{نقل‌قول:} \textit{«بی‌عدالتی در هر جا، تهدیدی برای عدالت در همه جاست.»}
\end{personbox}

\begin{personbox}[مالکوم ایکس (۱۹۲۵-۱۹۶۵)]

\textbf{نام اصلی:} مالکوم لیتل

\textbf{ملیت:} آمریکایی

\textbf{گرایش:} ناسیونالیسم سیاه، اسلام (ملت اسلام، سپس سنی)

\mediumspace

\textbf{تفاوت با کینگ:}
\begin{itemize}[nosep]
    \item نقد یکپارچگی با سفیدها: «جدایی، نه یکپارچگی»
    \item دفاع از حق دفاع مسلحانه
    \item نقد «خواب رویای آمریکایی»
    \item در سال‌های آخر: تحول به سمت سوسیالیسم و ضد نژادپرستی جهانی
\end{itemize}

\mediumspace

\textbf{نقل‌قول:} \textit{«ما برابری نمی‌خواهیم؛ می‌خواهیم به رسمیت شناخته شویم.»}
\end{personbox}

\begin{defbox}[تعریف: پلنگان سیاه]
\textbf{حزب پلنگ سیاه برای دفاع شخصی (۱۹۶۶):}
\begin{itemize}[nosep]
    \item بنیان‌گذاران: هیوی نیوتن، بابی سیل، الدریج کلیور
    \item ترکیب ناسیونالیسم سیاه با مارکسیسم-لنینیسم
    \item گشت‌زنی مسلحانه علیه خشونت پلیس
    \item برنامه‌های اجتماعی: صبحانه رایگان برای کودکان
    \item سرکوب شدید FBI (COINTELPRO)
\end{itemize}
\end{defbox}

\subsection{فمینیسم موج دوم و سوم}

\begin{figure}[htbp]
\centering
\begin{tikzpicture}[scale=0.85]

% عنوان
\node[font=\large\bfseries, text=DarkGray] at (0, 6.5) {\fa{سه موج فمینیسم}};

% موج اول
\fill[PurpleLight, opacity=0.7, rounded corners=8pt] (-6.5, 2) rectangle (-2, 5.5);
\draw[line width=2pt, color=PurpleAccent, rounded corners=8pt] (-6.5, 2) rectangle (-2, 5.5);
\node[text=PurpleAccent, font=\bfseries] at (-4.25, 5) {\fa{موج اول}};
\node[text=DarkGray, font=\small] at (-4.25, 4.4) {\fa{۱۸۵۰-۱۹۲۰}};
\node[text=DarkGray, font=\scriptsize, text width=3.5cm, align=center] at (-4.25, 3.3) {
    \fa{حق رأی}\\
    \fa{حق مالکیت}\\
    \fa{سوفراژت‌ها}
};

% موج دوم
\fill[PurpleAccent, opacity=0.3, rounded corners=8pt] (-1.5, 2) rectangle (2.5, 5.5);
\draw[line width=2pt, color=PurpleAccent, rounded corners=8pt] (-1.5, 2) rectangle (2.5, 5.5);
\node[text=PurpleAccent, font=\bfseries] at (0.5, 5) {\fa{موج دوم}};
\node[text=DarkGray, font=\small] at (0.5, 4.4) {\fa{۱۹۶۰-۱۹۸۰}};
\node[text=DarkGray, font=\scriptsize, text width=3.5cm, align=center] at (0.5, 3.3) {
    \fa{برابری شغلی}\\
    \fa{حقوق باروری}\\
    \fa{«شخصی سیاسی است»}
};

% موج سوم
\fill[PurpleAccent, opacity=0.5, rounded corners=8pt] (3, 2) rectangle (7, 5.5);
\draw[line width=2pt, color=PurpleAccent, rounded corners=8pt] (3, 2) rectangle (7, 5.5);
\node[text=PurpleAccent, font=\bfseries] at (5, 5) {\fa{موج سوم}};
\node[text=DarkGray, font=\small] at (5, 4.4) {\fa{۱۹۹۰-اکنون}};
\node[text=DarkGray, font=\scriptsize, text width=3.5cm, align=center] at (5, 3.3) {
    \fa{تقاطع‌گرایی}\\
    \fa{تنوع زنانه‌گی}\\
    \fa{نقد جنسیت دوگانه}
};

\end{tikzpicture}
\caption{سه موج جنبش فمینیستی}
\label{fig:feminism-waves}
\end{figure}

\begin{personbox}[سیمون دو بووار (۱۹۰۸-۱۹۸۶)]

\textbf{ملیت:} فرانسوی

\textbf{اثر کلیدی:} \book{جنس دوم} (۱۹۴۹)

\textbf{جمله معروف:} «زن زاده نمی‌شود، زن می‌شود.»

\mediumspace

\person{دوبووار} پایه‌گذار فمینیسم اگزیستانسیالیستی بود:
\begin{itemize}[nosep]
    \item زن به عنوان «دیگری» ساخته شده است
    \item زنانگی ساخته اجتماعی است، نه طبیعی
    \item رهایی زنان نیازمند استقلال اقتصادی و جنسی است
\end{itemize}
\end{personbox}

\begin{defbox}[تعریف: «شخصی سیاسی است»]
شعار کلیدی فمینیسم موج دوم که استدلال می‌کرد:
\begin{itemize}[nosep]
    \item تقسیم عمومی/خصوصی خود سیاسی است
    \item خانواده، جنسیت، بدن—همه عرصه‌های قدرتند
    \item ستم بر زنان فقط در حوزه عمومی نیست
    \item زندگی روزمره، عرصه مبارزه است
\end{itemize}
\end{defbox}

\subsection{جنبش LGBTQ+}

\begin{historybox}[شورش استون‌وال (۱۹۶۹)]
در ۲۸ ژوئن ۱۹۶۹، پلیس به بار استون‌وال این در نیویورک حمله کرد. این بار، همجنس‌گرایان مقاومت کردند:
\begin{itemize}[nosep]
    \item شش شب درگیری و اعتراض
    \item تولد جنبش حقوق همجنس‌گرایان مدرن
    \item تأسیس سازمان‌های فعال
    \item Pride Parade سالانه از همین‌جا آمد
\end{itemize}

\textbf{از «همجنس‌گرایی بیماری است» تا ازدواج برابر:} مسیری طولانی طی شد.
\end{historybox}

\begin{personbox}[جودیت باتلر (متولد ۱۹۵۶)]

\textbf{ملیت:} آمریکایی

\textbf{اثر کلیدی:} \book{دردسر جنسیت} (۱۹۹۰)

\textbf{ایده محوری:} جنسیت به مثابه عملکرد (Performance)

\mediumspace

\person{باتلر} بنیان‌گذار نظریه کوئیر است:
\begin{itemize}[nosep]
    \item جنسیت (Gender) ساخته فرهنگی است، نه طبیعی
    \item ما جنسیت را «اجرا» می‌کنیم، نه «هستیم»
    \item دوگانه مرد/زن خود سازه قدرت است
    \item هویت‌های نامتعارف (کوئیر) این نظم را به چالش می‌کشند
\end{itemize}
\end{personbox}

% ══════════════════════════════════════════════════════════════════════════════
\section{جنبش زیست‌محیطی}
% ══════════════════════════════════════════════════════════════════════════════

\begin{defbox}[تعریف: اکولوژی سیاسی]
جنبش زیست‌محیطی در دهه ۱۹۷۰ سیاسی شد:
\begin{itemize}[nosep]
    \item بحران محیط زیست ناشی از سیستم اقتصادی-سیاسی است
    \item رشد نامحدود در سیاره محدود ممکن نیست
    \item عدالت محیط‌زیستی: فقرا بیشتر آسیب می‌بینند
    \item نیاز به تغییر ساختاری، نه فقط «سبز زندگی کردن»
\end{itemize}
\end{defbox}

\begin{defbox}[تعریف: اکوسوسیالیسم]
ترکیب سوسیالیسم و اکولوژی:
\begin{itemize}[nosep]
    \item سرمایه‌داری ذاتاً ضد محیط زیست است (رشد اجباری)
    \item سوسیالیسم سنتی هم صنعت‌محور بود—باید اصلاح شود
    \item تولید باید برای نیاز باشد، نه سود
    \item دموکراسی اقتصادی برای تصمیم‌گیری درباره رابطه با طبیعت
\end{itemize}
\textbf{چهره‌ها:} ژوئل کوول، جان بلامی فاستر
\end{defbox}

\begin{defbox}[تعریف: رشدزدایی (Degrowth)]
نقد رادیکال‌تر که می‌گوید:
\begin{itemize}[nosep]
    \item «رشد سبز» توهم است
    \item باید مصرف و تولید را کاهش دهیم
    \item GDP معیار خوبی نیست
    \item زندگی ساده‌تر و بهتر ممکن است
\end{itemize}
\end{defbox}

% ══════════════════════════════════════════════════════════════════════════════
\section{آلترگلوبالیسم}
% ══════════════════════════════════════════════════════════════════════════════

\begin{historybox}[سیاتل ۱۹۹۹: تولد جنبش ضد جهانی‌سازی]
در نوامبر ۱۹۹۹، ده‌ها هزار نفر در سیاتل اجلاس WTO را مختل کردند:
\begin{itemize}[nosep]
    \item ائتلاف رنگارنگ: اتحادیه‌ها، محیط‌زیست‌گرایان، آنارشیست‌ها، کشاورزان
    \item شعار: «جهان دیگری ممکن است»
    \item روش: اقدام مستقیم، بلوکاد، اعتراض خلاقانه
    \item رسانه‌های جایگزین: ایندی‌مدیا
\end{itemize}

\textbf{ادامه:} پراگ ۲۰۰۰، جنوا ۲۰۰۱ (یک کشته)، کانکون ۲۰۰۳

\textbf{فوروم جهانی اجتماعی:} جایگزین داووس، از ۲۰۰۱ در پورتوآلگره برزیل
\end{historybox}

\begin{personbox}[نائومی کلاین (متولد ۱۹۷۰)]

\textbf{ملیت:} کانادایی

\textbf{آثار کلیدی:} \book{بدون لوگو} (۱۹۹۹)، \book{دکترین شوک} (۲۰۰۷)، \book{این همه چیز را تغییر می‌دهد} (۲۰۱۴)

\mediumspace

\person{کلاین} صدای برجسته جنبش آلترگلوبالیسم است:
\begin{itemize}[nosep]
    \item نقد برندها و فرهنگ مصرفی
    \item «سرمایه‌داری فاجعه»: استفاده از بحران برای خصوصی‌سازی
    \item پیوند بحران اقلیمی و نقد سرمایه‌داری
\end{itemize}
\end{personbox}

% ══════════════════════════════════════════════════════════════════════════════
\section{چپ پست‌مدرن و پسااستعماری}
% ══════════════════════════════════════════════════════════════════════════════

\begin{personbox}[میشل فوکو (۱۹۲۶-۱۹۸۴)]

\textbf{ملیت:} فرانسوی

\textbf{آثار کلیدی:} \book{نظم گفتار}، \book{مراقبت و تنبیه}، \book{تاریخ جنسیت}

\textbf{ایده محوری:} قدرت/دانش

\mediumspace

\person{فوکو} تحلیل قدرت را دگرگون کرد:
\begin{itemize}[nosep]
    \item قدرت فقط در دولت نیست—همه‌جا پراکنده است
    \item دانش و قدرت به هم پیوسته‌اند
    \item «انضباط» از طریق نهادها (زندان، بیمارستان، مدرسه)
    \item «زیست‌سیاست»: مدیریت جمعیت‌ها
    \item نقد «سوژه» و هویت ثابت
\end{itemize}
\end{personbox}

\begin{personbox}[ادوارد سعید (۱۹۳۵-۲۰۰۳)]

\textbf{ملیت:} فلسطینی-آمریکایی

\textbf{اثر کلیدی:} \book{شرق‌شناسی} (۱۹۷۸)

\textbf{ایده محوری:} شرق‌شناسی به مثابه گفتمان قدرت

\mediumspace

\person{سعید} بنیان‌گذار مطالعات پسااستعماری است:
\begin{itemize}[nosep]
    \item «شرق» ساخته غرب است
    \item دانش درباره شرق، ابزار سلطه استعماری بود
    \item بازنمایی‌ها (ادبیات، هنر، علم) سیاسی‌اند
    \item نقد امپریالیسم فرهنگی
\end{itemize}
\end{personbox}

\begin{debatebox}[مارکسیسم در برابر پست‌مدرنیسم]
\textbf{نقد مارکسیست‌ها به پست‌مدرنیسم:}
\begin{itemize}[nosep]
    \item نسبی‌گرایی: اگر همه چیز گفتمان است، حقیقت کجاست؟
    \item غفلت از اقتصاد: تمرکز بر فرهنگ، فراموشی طبقه
    \item بی‌برنامگی: نقد هست، راه‌حل نیست
    \item نخبه‌گرایی: زبان پیچیده، دور از توده‌ها
\end{itemize}

\textbf{پاسخ پست‌مدرن‌ها:}
\begin{itemize}[nosep]
    \item مارکسیسم خود «کلان‌روایتی» سرکوبگر است
    \item قدرت فقط اقتصادی نیست
    \item ما ابزارهای تحلیل می‌دهیم، نه نسخه
    \item چپ سنتی نژادپرستی و جنسیت‌گرایی خودش را نمی‌بیند
\end{itemize}
\end{debatebox}

% ══════════════════════════════════════════════════════════════════════════════
\section{چپ معاصر}
% ══════════════════════════════════════════════════════════════════════════════

\begin{figure}[htbp]
\centering
\begin{tikzpicture}[scale=0.85]

% عنوان
\node[font=\large\bfseries, text=DarkGray] at (0, 7) {\fa{چپ معاصر: بازگشت یا تحول؟}};

% تایم‌لاین
\draw[line width=4pt, MediumGray] (-7, 5) -- (7, 5);

% بحران ۲۰۰۸
\fill[OrangeAccent] (-5.5, 5) circle (10pt);
\node[above, text=OrangeAccent, font=\scriptsize\bfseries] at (-5.5, 5.5) {\fa{۲۰۰۸}};
\node[below, text width=2cm, align=center, text=DarkGray, font=\tiny] at (-5.5, 4.3) {
    \fa{بحران مالی}\\
    \fa{شکست نئولیبرالیسم؟}
};

% اشغال
\fill[LeftRed] (-3, 5) circle (10pt);
\node[above, text=LeftRed, font=\scriptsize\bfseries] at (-3, 5.5) {\fa{۲۰۱۱}};
\node[below, text width=2cm, align=center, text=DarkGray, font=\tiny] at (-3, 4.3) {
    \fa{اشغال وال‌استریت}\\
    \fa{بهار عربی}
};

% سیریزا/پودموس
\fill[CenterGreen] (-0.5, 5) circle (10pt);
\node[above, text=CenterGreen, font=\scriptsize\bfseries] at (-0.5, 5.5) {\fa{۲۰۱۵}};
\node[below, text width=2.3cm, align=center, text=DarkGray, font=\tiny] at (-0.5, 4.3) {
    \fa{سیریزا، پودموس}\\
    \fa{چپ رادیکال اروپا}
};

% سندرز/کوربین
\fill[LeftRedLight] (2, 5) circle (10pt);
\node[above, text=LeftRedLight, font=\scriptsize\bfseries] at (2, 5.5) {\fa{۱۵-۲۰}};
\node[below, text width=2cm, align=center, text=DarkGray, font=\tiny] at (2, 4.3) {
    \fa{سندرز، کوربین}\\
    \fa{بازگشت سوسیالیسم}
};

% کرونا
\fill[PurpleAccent] (4.5, 5) circle (10pt);
\node[above, text=PurpleAccent, font=\scriptsize\bfseries] at (4.5, 5.5) {\fa{۲۰۲۰}};
\node[below, text width=2cm, align=center, text=DarkGray, font=\tiny] at (4.5, 4.3) {
    \fa{کرونا}\\
    \fa{بازگشت دولت؟}
};

% BLM
\fill[LeftRedDark] (6.5, 5) circle (10pt);
\node[above, text=LeftRedDark, font=\scriptsize\bfseries] at (6.5, 5.5) {\fa{۲۰۲۰}};
\node[below, text width=2cm, align=center, text=DarkGray, font=\tiny] at (6.5, 4.3) {
    \fa{جان سیاهان مهم است}
};

\end{tikzpicture}
\caption{چپ در قرن بیست‌ویکم}
\label{fig:contemporary-left}
\end{figure}

\begin{personbox}[برنی سندرز (متولد ۱۹۴۱)]

\textbf{ملیت:} آمریکایی

\textbf{سمت:} سناتور ورمونت

\textbf{کارزار:} انتخابات مقدماتی دموکرات ۲۰۱۶ و ۲۰۲۰

\mediumspace

\person{سندرز} واژه «سوسیالیسم» را به گفتمان آمریکایی بازگرداند:
\begin{itemize}[nosep]
    \item بهداشت عمومی برای همه (Medicare for All)
    \item دانشگاه رایگان
    \item افزایش حداقل دستمزد به ۱۵ دلار
    \item مالیات بر میلیاردرها
    \item «سوسیالیسم دموکراتیک» به سبک اسکاندیناوی
\end{itemize}

\mediumspace

\textbf{اهمیت:} نسل جدیدی از جوانان آمریکایی به «سوسیالیسم» علاقه‌مند شدند.
\end{personbox}

\subsection{اختلافات درون چپ معاصر}

\begin{debatebox}[طبقه در برابر هویت]
\textbf{اولویت طبقه:}
\begin{itemize}[nosep]
    \item نابرابری اقتصادی ریشه همه مشکلات است
    \item تمرکز بر هویت، طبقه کارگر را پراکنده می‌کند
    \item «سیاست هویت» کارگران سفید را به سمت راست می‌راند
\end{itemize}

\textbf{تقاطع‌گرایی:}
\begin{itemize}[nosep]
    \item ستم‌ها (طبقه، نژاد، جنسیت) به هم پیوسته‌اند
    \item نمی‌توان یکی را به دیگری تقلیل داد
    \item «طبقه‌گرایی» خود می‌تواند سفیدمحور باشد
\end{itemize}

\textbf{این بحث، چپ معاصر را به دو اردوگاه تقسیم کرده است.}
\end{debatebox}

% ══════════════════════════════════════════════════════════════════════════════
%                         خلاصه و پرسش‌ها
% ══════════════════════════════════════════════════════════════════════════════

\newpage

\begin{summarybox}

\textbf{۱. چپ نو:} در دهه ۱۹۶۰ علیه هم سرمایه‌داری و هم چپ سنتی برآمد. تمرکز بر فرهنگ، آزادی، هویت.

\textbf{۲. ریشه‌های فکری:} مکتب فرانکفورت، به‌ویژه مارکوزه. نقد تمدن صنعتی، صنعت فرهنگ، سرکوب مازاد.

\textbf{۳. مه ۱۹۶۸:} لحظه نمادین—پاریس، پراگ، مکزیکو، آمریکا. «واقع‌بین باش، ناممکن را بخواه!»

\textbf{۴. جنبش‌های هویتی:} حقوق مدنی سیاهان (کینگ، مالکوم)، فمینیسم موج دوم و سوم، LGBTQ+.

\textbf{۵. جنبش زیست‌محیطی:} از حفاظت‌گرایی به اکوسوسیالیسم و رشدزدایی.

\textbf{۶. آلترگلوبالیسم:} سیاتل ۱۹۹۹، «جهان دیگری ممکن است».

\textbf{۷. چپ پست‌مدرن:} فوکو (قدرت پراکنده)، سعید (شرق‌شناسی)، باتلر (جنسیت). مناقشه با مارکسیسم.

\textbf{۸. امروز:} اشغال وال‌استریت، سندرز، کوربین، سیریزا. بحث طبقه/هویت ادامه دارد.

\end{summarybox}

\vspace{20pt}

\begin{questionbox}

\begin{enumerate}[nosep, rightmargin=1cm]
    \item آیا مه ۱۹۶۸ «شکست موفق» بود؟ اگر قدرت را نگرفت، چه چیزی را تغییر داد؟

    \item آیا تمرکز چپ نو بر فرهنگ و هویت، باعث غفلت از طبقه و اقتصاد شد؟

    \item نقد پست‌مدرنیست‌ها به «کلان‌روایت‌ها» شامل خود سوسیالیسم هم می‌شود. آیا این نقد موجه است؟

    \item آیا «سوسیالیسم دموکراتیک» سندرز واقعاً سوسیالیسم است، یا سوسیال‌دموکراسی اسکاندیناوی؟
\end{enumerate}

\end{questionbox}

% ══════════════════════════════════════════════════════════════════════════════
%                         پایان فصل هشتم
% ══════════════════════════════════════════════════════════════════════════════